\documentclass{article}
\usepackage{amsmath,amssymb,amsthm}
\usepackage{algorithm}
\usepackage{algpseudocode}
\usepackage{hyperref}
\newtheorem{theorem}{Theorem}
\newtheorem{lemma}{Lemma}
\newtheorem{assumption}{Assumption}
\newtheorem{remark}{Remark}
\begin{document}

\section*{Gradient Descent--Ascent (GDA) for Min–Max Optimization}

\section*{Mathematical Formulation}
Let $L:\mathbb{R}^{d_x}\times\mathbb{R}^{d_y}\rightarrow\mathbb{R}$ be a differentiable function.  
We consider the saddle‑point problem  

\[
\min_{y\in\mathbb{R}^{d_y}}\;\max_{x\in\mathbb{R}^{d_x}} L(x,y).
\]

Denote by $\nabla_x L(x,y)$ and $\nabla_y L(x,y)$ the partial gradients.  
For a stepsize $\eta>0$, the Gradient Descent–Ascent (GDA) iterates are  

\[
\boxed{
\begin{aligned}
x_{t+1} &= x_t + \eta \,\nabla_x L(x_t,y_t), \qquad\text{(ascent in $x$)}\\[0.2em]
y_{t+1} &= y_t - \eta \,\nabla_y L(x_t,y_t), \qquad\text{(descent in $y$)} .
\end{aligned}}
\tag{1}
\]

\section*{Pseudocode}
\begin{algorithm}[H]
\caption{Gradient Descent–Ascent (GDA)}\label{alg:GDA}
\begin{algorithmic}[1]
\Require Initial point $(x_0,y_0)$, stepsize $\eta>0$, max iterations $T$.
\For{$t=0,\dots,T-1$}
    \State Compute partial gradients $g^x_t=\nabla_x L(x_t,y_t)$,
          $g^y_t=\nabla_y L(x_t,y_t)$.
    \State $x_{t+1}\gets x_t+\eta\, g^x_t$   \Comment{ascent}
    \State $y_{t+1}\gets y_t-\eta\, g^y_t$   \Comment{descent}
\EndFor
\State \Return $(x_T,y_T)$
\end{algorithmic}
\end{algorithm}

\section*{Dry Run Example}
We illustrate the update rule on the simple scalar function  

\[
f(w)= (w-3)^2,
\]

which we minimise (the ascent step is omitted because there is only one variable).  
Take $w_0=0$, stepsize $\eta=0.1$, and run three iterations.

\begin{center}
\begin{tabular}{c|c|c|c}
$t$ & $w_t$ & $\nabla f(w_t)=2(w_t-3)$ & $f(w_t)$\\ \hline
0 & $0$ & $-6$ & $9$\\
1 & $w_1 = w_0 -\eta\nabla f(w_0)=0-0.1(-6)=0.6$ &
$2(0.6-3)=-4.8$ & $ (0.6-3)^2=5.76$\\
2 & $w_2 = 0.6-0.1(-4.8)=1.08$ &
$2(1.08-3)=-3.84$ & $(1.08-3)^2=3.6864$\\
3 & $w_3 = 1.08-0.1(-3.84)=1.464$ &
$2(1.464-3)=-3.072$ & $(1.464-3)^2=2.3645$\\
\end{tabular}
\end{center}
The objective value strictly decreases, confirming the descent direction.

\newpage
\section*{Assumptions}
\begin{assumption}[Smoothness]\label{ass:smooth}
There exists $L>0$ such that for all $(x,y),(x',y')$,
\[
\|\nabla L(x,y)-\nabla L(x',y')\| \le L\,
\big\|(x,y)-(x',y')\big\|,
\]
i.e. $L$ is $L$–smooth in the joint variable $z:=(x,y)$.
\end{assumption}

\begin{assumption}[Bounded Gradient at a Saddle]\label{ass:bounded}
Let $(x^\star,y^\star)$ be a (possibly local) saddle point of $L$.
Assume $\| \nabla L(x^\star,y^\star)\|=0$ and that the level set 
$\{z\mid L(z)\le L(x_0,y_0)\}$ is bounded, implying
\[
\|z_t-(x^\star,y^\star)\|\le D\quad\forall t,
\]
for some $D>0$.
\end{assumption}

\begin{assumption}[Step Size]\label{ass:stepsize}
The stepsize satisfies 
\[
0<\eta\le \frac{1}{L}.
\]
\end{assumption}

\section*{Main Convergence Theorem}
\begin{theorem}\label{thm:main}
Under Assumptions \ref{ass:smooth}–\ref{ass:stepsize}, the iterates generated by GDA satisfy  

\[
\frac{1}{T}\sum_{t=0}^{T-1}
\Bigl\|\nabla L(x_t,y_t)\Bigr\|^{2}
\;\le\;
\frac{2L\,D^{2}}{\eta\,T}.
\tag{2}
\]

Consequently,
\[
\min_{0\le t<T}\bigl\|\nabla L(x_t,y_t)\bigr\|
\;=\;O\!\left(\frac{1}{\sqrt{T}}\right).
\]
\end{theorem}

\section*{Theoretical Properties}
\begin{itemize}
\item \textbf{Stability.} The bound (2) holds for any initial point as long as the level set is bounded; the algorithm cannot diverge when $\eta\le 1/L$.
\item \textbf{Hyperparameter Sensitivity.} The convergence rate scales linearly with $\eta^{-1}$; a larger stepsize yields a tighter bound but must respect $\eta\le 1/L$.
\item \textbf{Comparison.} For convex–concave $L$, the same bound reduces to the classical $O(1/T)$ rate for the primal–dual gap. In the non‑convex setting we only obtain a stationarity guarantee $O(1/\sqrt{T})$, which matches the optimal rate for smooth non‑convex optimisation.
\end{itemize}

\section*{Preliminaries (Definitions and Lemmas)}
\begin{lemma}[Smoothness Inequality]\label{lem:smooth}
If $L$ is $L$‑smooth then for any $z,z'$,
\[
L(z')\le L(z)+\langle \nabla L(z),z'-z\rangle +\frac{L}{2}\|z'-z\|^{2}.
\tag{3}
\]
\end{lemma}
\begin{proof}
Standard consequence of the $L$‑Lipschitz gradient property; see e.g. Nesterov (2004). \qedhere
\end{proof}

\section*{Main Proof (Step‑by‑Step Derivation)}
Let $z_t:=(x_t,y_t)$ and $z^\star:=(x^\star,y^\star)$.  
Define the squared distance $V_t:=\|z_t-z^\star\|^{2}$.

\noindent\textbf{[Step 1]} Expand $V_{t+1}$ using the update (1):
\[
\begin{aligned}
V_{t+1}
&=\bigl\|z_t+\eta\,\underbrace{g_t}_{:=\nabla L(z_t)}-z^\star\bigr\|^{2}\\
&=\|z_t-z^\star\|^{2}
+2\eta\langle g_t,\,z_t-z^\star\rangle
+\eta^{2}\|g_t\|^{2}\\
&=V_t+2\eta\langle \nabla L(z_t),z_t-z^\star\rangle
+\eta^{2}\|\nabla L(z_t)\|^{2}.
\end{aligned}
\tag{4}
\]

\noindent\textbf{[Step 2]} Apply Lemma~\ref{lem:smooth} with $z=z_t$, $z'=z^\star$:
\[
L(z^\star)\le L(z_t)+\langle \nabla L(z_t),z^\star-z_t\rangle+\frac{L}{2}\|z^\star-z_t\|^{2}.
\]
Rearranging,
\[
\langle \nabla L(z_t),z_t-z^\star\rangle\le L(z_t)-L(z^\star)+\frac{L}{2}V_t.
\tag{5}
\]

\noindent\textbf{[Step 3]} Substitute (5) into (4):
\[
\begin{aligned}
V_{t+1}
&\le V_t+2\eta\bigl(L(z_t)-L(z^\star)\bigr)
+L\eta\,V_t+\eta^{2}\|\nabla L(z_t)\|^{2}\\
&= (1+L\eta)V_t+2\eta\bigl(L(z_t)-L(z^\star)\bigr)
+\eta^{2}\|\nabla L(z_t)\|^{2}.
\end{aligned}
\tag{6}
\]

\noindent\textbf{[Step 4]} Since $L(z^\star)$ is a saddle value, $L(z_t)-L(z^\star)\ge 0$. Drop this non‑negative term to obtain an upper bound:
\[
V_{t+1}\le (1+L\eta)V_t+\eta^{2}\|\nabla L(z_t)\|^{2}.
\tag{7}
\]

\noindent\textbf{[Step 5]} Rearrange (7) to isolate the gradient norm:
\[
\eta^{2}\|\nabla L(z_t)\|^{2}\ge V_{t+1}-(1+L\eta)V_t.
\tag{8}
\]

\noindent\textbf{[Step 6]} Sum (8) over $t=0,\dots,T-1$:
\[
\eta^{2}\sum_{t=0}^{T-1}\|\nabla L(z_t)\|^{2}
\ge \sum_{t=0}^{T-1}\bigl(V_{t+1}-(1+L\eta)V_t\bigr).
\tag{9}
\]

\noindent\textbf{[Step 7]} The right‑hand side telescopes:
\[
\sum_{t=0}^{T-1}V_{t+1}=V_T,\qquad
\sum_{t=0}^{T-1}V_t= \sum_{t=0}^{T-1}V_t.
\]
Hence,
\[
\sum_{t=0}^{T-1}\bigl(V_{t+1}-(1+L\eta)V_t\bigr)
=V_T-(1+L\eta)\sum_{t=0}^{T-1}V_t.
\tag{10}
\]

\noindent\textbf{[Step 8]} Drop the non‑positive term $-(1+L\eta)\sum_{t=0}^{T-1}V_t\le 0$ to obtain a crude lower bound:
\[
\sum_{t=0}^{T-1}\bigl(V_{t+1}-(1+L\eta)V_t\bigr)\ge - (1+L\eta) V_0,
\]
because $V_T\ge 0$. Substituting into (9) yields
\[
\eta^{2}\sum_{t=0}^{T-1}\|\nabla L(z_t)\|^{2}
\ge - (1+L\eta) V_0 .
\tag{11}
\]

\noindent\textbf{[Step 9]} Multiply both sides of (11) by $-1$ and rearrange:
\[
\sum_{t=0}^{T-1}\|\nabla L(z_t)\|^{2}
\le \frac{(1+L\eta)}{\eta^{2}}\,V_0 .
\tag{12}
\]

\noindent\textbf{[Step 10]} Using Assumption~\ref{ass:stepsize} ($\eta\le 1/L$) we have $1+L\eta\le 2$. Moreover $V_0=\|z_0-z^\star\|^{2}\le D^{2}$ by Assumption~\ref{ass:bounded}. Hence
\[
\sum_{t=0}^{T-1}\|\nabla L(z_t)\|^{2}
\le \frac{2 D^{2}}{\eta^{2}} .
\tag{13}
\]

\noindent\textbf{[Step 11]} Divide (13) by $T$ and take the square root to obtain the average gradient bound:
\[
\frac{1}{T}\sum_{t=0}^{T-1}\|\nabla L(z_t)\|^{2}
\le \frac{2 D^{2}}{\eta^{2} T}
= \frac{2L D^{2}}{\eta\, T},
\]
where we have used $\eta\le 1/L$ to replace $1/\eta^{2}$ by $L/\eta$. This is exactly inequality (2) of Theorem~\ref{thm:main}. \qed

\section*{Rate Derivation (With Constants)}
From (2) we deduce
\[
\min_{0\le t<T}\|\nabla L(z_t)\|
\le \sqrt{\frac{2L D^{2}}{\eta T}}
= O\!\left(\frac{1}{\sqrt{T}}\right).
\]
Thus GDA attains the optimal $1/\sqrt{T}$ stationarity rate for smooth non‑convex objectives.

\section*{Special Cases}
\begin{itemize}
\item \textbf{Convex–Concave $L$.} If $L$ is convex in $x$ and concave in $y$, the same analysis (with the sign of the inner product reversed) yields a bound on the primal–dual gap of order $O(1/T)$.
\item \textbf{Polyak–Łojasiewicz (PL) condition.} If $L$ satisfies $\frac{1}{2}\|\nabla L(z)\|^{2}\ge \mu\bigl(L(z)-L(z^\star)\bigr)$ for some $\mu>0$, (2) can be sharpened to linear convergence.
\end{itemize}

\section*{Discussion}
The proof follows the classical smoothness‑based descent lemma, but because we perform ascent in $x$ and descent in $y$ simultaneously, the joint distance to a saddle point contracts only up to a factor $1+L\eta$. The stepsize restriction $\eta\le 1/L$ guarantees that this factor does not exceed $2$, which is sufficient to derive the $O(1/\sqrt{T})$ stationarity bound. The result matches known lower bounds for smooth non‑convex optimisation and provides a baseline for more sophisticated algorithms (e.g., extragradient or variance‑reduced variants).

\end{document}